% basic nastaveni
\documentclass[12pt]{article}
\setlength{\parindent}{0pt}
\setlength{\parskip}{1em}
\usepackage[utf8]{inputenc}
\usepackage[czech]{babel}
\usepackage[T1]{fontenc}
\usepackage{graphics}
\usepackage{caption}
\usepackage{hyperref}
\usepackage{dsfont}
\usepackage{color}
\usepackage{float}
\usepackage{enumitem}
\usepackage{graphicx}

% matematicke package
\usepackage{tikz-cd}
\usepackage{tikz}
\usetikzlibrary{babel}
\usepackage{amsmath}
\usepackage{amssymb}
\usepackage{amsfonts}
\usepackage{amssymb}
\usepackage{geometry}
\geometry{margin=2.5cm}
\usepackage{pgfplots}
\pgfplotsset{compat=1.18}
\usepackage{mathrsfs}
\usetikzlibrary{shapes.misc}
\usetikzlibrary{angles,quotes}

\title{Maturitní otázka číslo 10: Pravděpodobnost}
\author{}
\date{}

\begin{document}
\maketitle

\section*{Základní pojmy}

\begin{itemize}[itemsep = 2pt]
    \item \textbf{Náhodný pokus} je proces, nebo pozorování, jehož výsledek nelze předem určit. Je prováděn
    za stále stejných, předem stanovených podmínek (hod kostkou)
    \item \textbf{Množina všech možných výsledků ($\Omega$)} je neprázdná množina všech výsledků, které mohou při náhodném pokusu nastat
    \item \textbf{Jev} je podmnožina možných výsledků náhodného pokusu, o které lze následně říci, jestli nastala, či nikoliv
    \item \textbf{Jistý jev} je jev, který nastane při každém pokusu
    \item \textbf{Nemožný jev} je jev, který nikdy nenastane ($\Omega = \{\emptyset\}$)
\end{itemize}

\subsection*{Množinové operace s jevy}
\begin{itemize}[itemsep = 2pt]
    \item Sjednocení ($A \cup B$): Nastává jev, pokud nastane alespoň jeden z jevů  $A$ nebo $B$.
    \item Průnik ($A \cap B$): Nastává jev, pokud nastanou jevy $A$ a $B$ současně
    \item Doplňek / opačný jev ($A'$): Nastává jev, že jev $A$ nenastane.
    \item $B \subset A$: $B$ je podjev jevu $A$ 
\end{itemize}

\section*{Klasická definice pravděpodobnosti}
Pokud má všech $m$ možných výsledků náhodného pokusu stejnou pravděpodobnost vzniku a jev $A$ je tvořen $m(A)$ příznivými výsledky, pak  pravděpodobnost, že nastane jev $A$ je: 
\[P(A) = \frac{m(A)}{m}\]

\subsection*{Operace s pravděpodobnostmi}
\subsubsection*{Pravdpodobnost opačného jevu}
Jelikož jev $A$ a jeho doplněk $A'$ tvoří celou množinu $\Omega$, pak platí:
\[P(A') = 1-P(A)\]

\subsubsection*{Pravděpodobnost sjednocení jevů}
Pokud jsou jevy $A$ a $B$ neslučitelné (nemohou se stát zároveň, $A \cap B = \emptyset$), pak platí:
\[P(A\cup B) = P(A) + P(B)\]

Pokud jsou jevy slučitelné ($A \cap B \neq \emptyset$), pak platí:
\[P(A\cup B) = P(A) + P(B) - P(A\cap B)\]

\section*{Nezávislé jevy}
Dva jevy $A$ a $B$ jsou nezávislé, pokud výskyt jednoho jevu neovliňuje pravděpodobnost výskytu jevu druhého.
 Např. hod mincí, pravděpodobnost jejich průniku je součinem jejich pravděpodobností (stanou se oba dva).
$P(A\cap B) = P(A) \cdot P(B)$.

\subsection*{Opakované nezávislé pokusy}
Jde o situaci, kdy provádíme týž pokus vícekrát za stejných podmínek a výsledky jendotlivých pokusů se navzájem neovliňují.

\subsubsection*{Binomické rozdělení pravděpodobnosti (Bernoulliho vzorec)}
Používá se pro výpočet pravděpodobnosti, kde při $n$ opakovaných nezávislých pokusech nastane jev $A$ právě $k$-krát. Jeho znění je následovné:
\[P_n(k) = \binom{n}{k}\cdot p^k \cdot q^{n-k}\]
$p$ je pravděpodobnost úspěchu v jednom dílčím pokusu a $q$ je pravděpodobnost neúspěchu v jednom pokusu ($q = 1-p$)




\end{document}